\section{Lecture 7: ANOVA, Model Comparison, and $R^2$}

\subsection{Decomposition of Variation}

For each observation,
\[
y_i - \bar{y} = (y_i - \hat{y}_i) + (\hat{y}_i - \bar{y}).
\]

Total deviation from the mean splits into a residual part and a model-explained part.

\subsection{Sum of Squares}

\begin{defn}{Total Sum of Squares}{def:sst}
Total variation in $y$:
\[
\SST = s_{yy} = \sum_{i=1}^{n}(y_i-\bar{y})^2.
\]
\end{defn}

\begin{defn}{Residual Sum of Squares}{def:rss}
Unexplained variation:
\[
\RSS=\sum_{i=1}^{n}(y_i-\hat{y}_i)^2.
\]
\end{defn}

\begin{defn}{Regression Sum of Squares}{def:ssreg}
Variation explained by the model:
\[
\SSreg=\sum_{i=1}^{n}(\hat{y}_i-\bar{y})^2.
\]
\end{defn}

\subsection{ANOVA Decomposition}

Total variation splits into explained and unexplained parts:
\[
\SST = \SSreg + \RSS.
\]

\textbf{Interpretation:} $\SSreg$ measures variation explained by the regression; $\RSS$ measures variation left in the residuals. If $\hat{y}_i \approx \bar{y}$ for all $i$, then $\SSreg \approx 0$ (model explains little). If many $\hat{y}_i$ differ strongly from $\bar{y}$, then $\SSreg$ is large (model explains substantial variation).

\subsection{Symmetric and Idempotent Matrices}

\begin{defn}{Symmetric Matrix}{def:symmetric}
A square matrix $A$ is symmetric if $A^\top = A$.
\end{defn}

\begin{defn}{Idempotent Matrix}{def:idempotent}
A square matrix $A$ is idempotent if $A^2 = A$.
\end{defn}

\subsection{Hat Matrix}

The regression projection matrix is
\[
H = \bX(\bX^\top \bX)^{-1}\bX^\top.
\]

The hat matrix satisfies $H^\top = H$ and $H^2 = H$; hence $H$ is symmetric and idempotent.

\textbf{Proof sketch:}
\[
H^\top = (\bX(\bX^\top \bX)^{-1}\bX^\top)^\top
= \bX(\bX^\top \bX)^{-1}\bX^\top
= H
\]
(using $(ABC)^\top = C^\top B^\top A^\top$ and $(A^\top)^{-1} = (A^{-1})^\top$).

For idempotence,
\[
H^2 = \bX(\bX^\top \bX)^{-1}\bX^\top \bX(\bX^\top \bX)^{-1}\bX^\top
= \bX(\bX^\top \bX)^{-1}(\bX^\top \bX)(\bX^\top \bX)^{-1}\bX^\top
= \bX(\bX^\top \bX)^{-1}\bX^\top = H.
\]

\subsection{Vector of Ones and Matrix of Ones}

Let
\[
\mathbf{1} =
\begin{pmatrix}
1\\
\vdots\\
1
\end{pmatrix}, \qquad
J = \mathbf{1}\mathbf{1}^{\top}.
\]

Then $J$ is an $n \times n$ matrix of ones:
\[
J =
\begin{pmatrix}
1 & \cdots & 1\\
\vdots & \ddots & \vdots\\
1 & \cdots & 1
\end{pmatrix}.
\]

\subsection{Properties of $J$}

\textbf{Dimension:} $n \times n$.

\textbf{Transpose:}
\[
J^{\top} = (\mathbf{1}\mathbf{1}^{\top})^{\top}
= (\mathbf{1}^{\top})^{\top}\mathbf{1}^{\top}
= \mathbf{1}\mathbf{1}^{\top}
= J.
\]
Thus $J$ is symmetric.

\textbf{Square:}
\[
J^2 = (\mathbf{1}\mathbf{1}^{\top})(\mathbf{1}\mathbf{1}^{\top})
= \mathbf{1}(\mathbf{1}^{\top}\mathbf{1})\mathbf{1}^{\top}.
\]
Since $\mathbf{1}^{\top}\mathbf{1}=n$,
\[
J^2 = n\mathbf{1}\mathbf{1}^{\top} = nJ.
\]
Hence $J$ is symmetric but \emph{not} idempotent.

\subsection{Scaled Ones Matrix}

Define $\frac{1}{n}J$. Then

\textbf{Transpose:}
\[
\left(\frac{1}{n}J\right)^{\top}=\frac{1}{n}J.
\]

\textbf{Square:}
\[
\left(\frac{1}{n}J\right)^2
= \frac{1}{n^2}J^2
= \frac{1}{n^2}(nJ)
= \frac{1}{n}J.
\]

Thus $\frac{1}{n}J$ is symmetric and idempotent.

\begin{notebox}
\textbf{Euclidean Norm:} For $x=(x_1,\dots,x_n)^{\top}$,
\[
\|x\|^2 = \sum_{i=1}^{n}x_i^2 = x^{\top}x.
\]
\end{notebox}

\subsection{Residual Maker Matrix}

Define the residual projection matrix
\[
M = I - H,
\]
where $H = \bX(\bX^{\top}\bX)^{-1}\bX^{\top}$.

\textbf{Symmetry:}
\[
(I-H)^{\top} = I^{\top} - H^{\top} = I - H.
\]

\textbf{Idempotence:}
\[
(I-H)^2 = (I-H)(I-H) = I - IH - HI + H^2.
\]
Using $IH=HI=H$ and $H^2=H$,
\[
(I-H)^2 = I - H - H + H = I - H.
\]

Thus $I-H$ is symmetric and idempotent.

\subsection{Mean Vector Representation}

Let $y=(y_1,\dots,y_n)^{\top}$ and $\mathbf{1}=(1,\dots,1)^{\top}$. The sample mean is
\[
\bar{y}=\frac{1}{n}\sum_{i=1}^n y_i.
\]

Define the mean vector
\[
\bar{y}\,\mathbf{1}=
\begin{pmatrix}
\bar{y}\\
\vdots\\
\bar{y}
\end{pmatrix}.
\]

Using the ones matrix $J=\mathbf{1}\mathbf{1}^{\top}$,
\[
\bar{y}\,\mathbf{1}=\frac{1}{n}Jy.
\]

\subsection{Total Sum of Squares (Matrix Form)}

Total variation:
\[
\SST=\sum_{i=1}^n (y_i-\bar{y})^2.
\]

Vector form:
\[
\SST=\|y-\bar{y}\mathbf{1}\|^2.
\]

Substitute $\bar{y}\mathbf{1}=\frac{1}{n}Jy$:
\[
\SST=\left\|y-\frac{1}{n}Jy\right\|^2
=\left\|(I-\tfrac{1}{n}J)y\right\|^2.
\]

Using $\|x\|^2=x^{\top}x$:
\[
\SST=((I-\tfrac{1}{n}J)y)^{\top}((I-\tfrac{1}{n}J)y)
= y^{\top}(I-\tfrac{1}{n}J)^{\top}(I-\tfrac{1}{n}J)y.
\]

Since $(I-\tfrac{1}{n}J)$ is symmetric and idempotent,
\[
\SST = y^{\top}(I-\tfrac{1}{n}J)y.
\]

\subsection{Residual Sum of Squares (Matrix Form)}

Residuals:
\[
\hat{y} = Hy.
\]

Residual vector:
\[
e = y-\hat{y} = y-Hy = (I-H)y.
\]

Residual sum of squares:
\[
\RSS=\sum_{i=1}^n (y_i-\hat{y}_i)^2
=\|y-\hat{y}\|^2.
\]

Matrix form:
\[
\RSS=(y-\hat{y})^{\top}(y-\hat{y})
=(y-Hy)^{\top}(y-Hy).
\]

\[
\RSS=((I-H)y)^{\top}((I-H)y)
= y^{\top}(I-H)^{\top}(I-H)y.
\]

Since $I-H$ is symmetric and idempotent,
\[
\RSS = y^{\top}(I-H)y.
\]

\subsection{Matrix $J$ and Its Properties}

Let $\mathbf{1}=(1,\dots,1)^\top \in \mathbb{R}^n$ and define
\[
J=\mathbf{1}\mathbf{1}^\top.
\]

\textbf{Symmetric:} A matrix $A$ is symmetric if $A^\top=A$.
\[
\left(\frac{1}{n}J\right)^\top
=\frac{1}{n}J^\top
=\frac{1}{n}(\mathbf{1}\mathbf{1}^\top)^\top
=\frac{1}{n}(\mathbf{1}^\top)^\top \mathbf{1}^\top
=\frac{1}{n}\mathbf{1}\mathbf{1}^\top
=\frac{1}{n}J.
\]

\textbf{Idempotent:} A matrix $A$ is idempotent if $A^2=A$.
\[
\left(\frac{1}{n}J\right)^2
=\frac{1}{n^2}JJ
=\frac{1}{n^2}\mathbf{1}\mathbf{1}^\top\mathbf{1}\mathbf{1}^\top
=\frac{1}{n^2}\mathbf{1}(\mathbf{1}^\top\mathbf{1})\mathbf{1}^\top.
\]

Since $\mathbf{1}^\top\mathbf{1}=n$,
\[
\left(\frac{1}{n}J\right)^2
=\frac{1}{n}\mathbf{1}\mathbf{1}^\top
=\frac{1}{n}J.
\]

\subsection{Hat Matrix Properties}

\[
H = \bX(\bX^\top \bX)^{-1}\bX^\top.
\]

\textbf{Properties:} $H$ is symmetric ($H^\top=H$) and idempotent ($H^2=H$).

\textbf{Fitted values:}
\[
\hat{y} = Hy.
\]

\textbf{Residuals:}
\[
e = y-\hat{y} = (I-H)y.
\]

\subsection{Interaction Between $H$ and $J$}

\begin{keyresult}[Interaction Property]
If the intercept is included in the regression,
\[
\frac{1}{n}HJ=\frac{1}{n}J, \qquad
\frac{1}{n}JH=\frac{1}{n}J.
\]
\end{keyresult}

If the intercept is included in the regression, then $\mathbf{1}\in\text{col}(\bX)$, so there exists $\alpha$ such that
\[
\mathbf{1}=\bX\alpha.
\]

Then
\[
H\mathbf{1}
= \bX(\bX^\top \bX)^{-1}\bX^\top\mathbf{1}
= \bX(\bX^\top \bX)^{-1}\bX^\top \bX\alpha
= \bX\alpha
= \mathbf{1}.
\]

Thus
\[
\frac{1}{n}HJ
=\frac{1}{n}H\mathbf{1}\mathbf{1}^\top
=\frac{1}{n}\mathbf{1}\mathbf{1}^\top
=\frac{1}{n}J.
\]

Similarly, $\frac{1}{n}JH=\frac{1}{n}J$.

\subsection{Regression Sum of Squares}

\[
\SSreg
=\sum_{i=1}^n(\hat{y}_i-\bar{y})^2
=\|\hat{y}-\bar{y}\mathbf{1}\|^2.
\]

Since
\[
\hat{y} = Hy, \qquad \bar{y}\mathbf{1}=\frac{1}{n}Jy,
\]

\[
\SSreg
=\|Hy-\tfrac{1}{n}Jy\|^2
=\|(H-\tfrac{1}{n}J)y\|^2.
\]

Using $\|Ax\|^2 = x^\top A^\top A x$,
\[
\SSreg
=((H-\tfrac{1}{n}J)y)^\top((H-\tfrac{1}{n}J)y)
= y^\top(H-\tfrac{1}{n}J)^\top(H-\tfrac{1}{n}J)y.
\]

Using symmetry and the identities above,
\[
\SSreg=y^\top(H-\tfrac{1}{n}J)y.
\]

\subsection{Matrix Forms of Sums of Squares}

\begin{keyresult}[Matrix Forms]
\[
\SST = y^\top\left(I-\frac{1}{n}J\right)y,
\]
\[
\RSS = y^\top(I-H)y,
\]
\[
\SSreg = y^\top\left(H-\frac{1}{n}J\right)y.
\]
\end{keyresult}

\subsection{ANOVA Decomposition}

\[
\SST = \RSS + \SSreg.
\]

\subsection{ANOVA Decomposition (Scalar Proof)}

\textbf{Goal:}
\[
\sum_{i=1}^{n}(y_i-\bar{y})^2
=
\sum_{i=1}^{n}(y_i-\hat{y}_i)^2
+
\sum_{i=1}^{n}(\hat{y}_i-\bar{y})^2.
\]

Write
\[
y_i-\bar{y}=(y_i-\hat{y}_i)+(\hat{y}_i-\bar{y}).
\]

Square and sum:
\[
\sum_{i=1}^{n}(y_i-\bar{y})^2
=
\sum_{i=1}^{n}\big[(y_i-\hat{y}_i)+(\hat{y}_i-\bar{y})\big]^2
\]

\[
=
\sum_{i=1}^{n}(y_i-\hat{y}_i)^2
+
\sum_{i=1}^{n}(\hat{y}_i-\bar{y})^2
+
2\sum_{i=1}^{n}(y_i-\hat{y}_i)(\hat{y}_i-\bar{y}).
\]

The cross-product term vanishes (see proof below), so
\[
\SST = \RSS + \SSreg.
\]

\subsection{Proof: $\sum_{i=1}^{n}e_i(\hat{y}_i-\bar{y}) = 0$}

Residuals: $e_i = y_i-\hat{y}_i$. Then
\[
\sum_{i=1}^{n} e_i(\hat{y}_i-\bar{y})
=
\sum_{i=1}^{n}e_i\hat{y}_i
-
\bar{y}\sum_{i=1}^{n}e_i.
\]

In least squares, $\sum_{i=1}^{n}e_i = 0$. For the first part,
\[
\sum_{i=1}^{n} e_i\hat{y}_i
=
\sum_{i=1}^{n} e_i
(\beta_0+\beta_1 x_{i1}+\cdots+\beta_p x_{ip})
\]

\[
=
\beta_0\sum_{i=1}^{n}e_i
+
\beta_1\sum_{i=1}^{n}e_i x_{i1}
+\cdots+
\beta_p\sum_{i=1}^{n}e_i x_{ip}.
\]

The least squares normal equations imply
\[
\sum_{i=1}^{n}e_i=0,
\qquad
\sum_{i=1}^{n}e_i x_{ij}=0 \quad (j=1,\dots,p).
\]

Therefore, $\sum_{i=1}^{n} e_i\hat{y}_i = 0$, so $\sum_{i=1}^{n}e_i(\hat{y}_i-\bar{y}) = 0$.

\subsection{Degrees of Freedom}

\[
\SST=\sum_{i=1}^{n}(y_i-\bar{y})^2
\quad\text{has } n-1 \text{ degrees of freedom}.
\]

\[
\RSS=\sum_{i=1}^{n}(y_i-\hat{y}_i)^2
\quad\text{has } n-p-1 \text{ degrees of freedom}.
\]

\[
\SSreg=\sum_{i=1}^{n}(\hat{y}_i-\bar{y})^2
\quad\text{has } p \text{ degrees of freedom}.
\]

\textbf{Check:} $n-1=(n-p-1)+p$.

\subsection{ANOVA F-Test Idea}

Previously, we tested $H_0:\beta_1=0$ using the $t$-statistic
\[
T=\frac{\bhat{\beta}_1-0}{\se(\bhat{\beta}_1)}.
\]

ANOVA provides another way to test whether the predictor contributes to the model.

\subsection{Comparing Two Models}

ANOVA compares two nested models.

\textbf{Model 1 (Reduced model):}
\[
Y_i=\beta_0+\varepsilon_i.
\]

\textbf{Model 2 (Full model):}
\[
Y_i=\beta_0+\beta_1 X_i+\varepsilon_i.
\]

\textbf{Hypothesis test:}
\[
H_0:\beta_1=0
\qquad
H_1:\beta_1\neq 0.
\]

Model 1 assumes no relationship between $X$ and $Y$, while Model 2 allows $X$ to explain variation in $Y$.

\subsection{Fitting the Models}

Fit both models and compute predicted values:
\[
\hat{y}_i^{(1)} \quad \text{from Model 1},
\qquad
\hat{y}_i^{(2)} \quad \text{from Model 2}.
\]

Residual Sum of Squares:
\[
\RSS_1=\sum_{i=1}^{n}(y_i-\hat{y}_i^{(1)})^2,
\qquad
\RSS_2=\sum_{i=1}^{n}(y_i-\hat{y}_i^{(2)})^2.
\]

\textbf{Interpretation:} $\RSS$ measures the amount of variation in $Y$ not explained by the model.

\subsubsection{Why $\RSS_1 \ge \RSS_2$}

For Model 1: $\hat{y}_i^{(1)}=\bar{y}$, so
\[
\RSS_1=\sum_{i=1}^{n}(y_i-\bar{y})^2=\SST.
\]

For Model 2: $\hat{y}_i^{(2)}=\hat{y}_i$, so
\[
\RSS_2=\sum_{i=1}^{n}(y_i-\hat{y}_i)^2=\RSS.
\]

Thus, $\RSS_1 \ge \RSS_2$.

\begin{caution}
\textbf{Is Model 2 Always Better?} Not necessarily. Model 2 always reduces $\RSS$ because it includes an extra parameter. The key question is whether the reduction is large enough to justify adding $\beta_1$.

If $\RSS_1 \gg \RSS_2$, Model 2 is better.
If $\RSS_1 \approx \RSS_2$, Model 1 is better.
\end{caution}

\subsection{Explained Variation}

The reduction in residual variation is $\RSS_1-\RSS_2$. Using the ANOVA identity:
\[
\SSreg=\RSS_1-\RSS_2=\SST-\RSS.
\]

Thus, the regression sum of squares measures how much variation in $Y$ is explained by the predictor.

\subsection{F Distribution}

Let
\[
F \sim F(df_1, df_2).
\]

Then $F$ follows an F-distribution with degrees of freedom $df_1>0$ and $df_2>0$.

\textbf{Key properties:}
\begin{itemize}
\item $F \ge 0$ (non-negative).
\item The distribution is right-skewed.
\item The exact shape depends on the degrees of freedom.
\end{itemize}

\subsubsection{Definition of the F Distribution}

If
\[
U \sim \chi^2_{df_1}, \qquad V \sim \chi^2_{df_2}
\]
and $U$ and $V$ are independent, then
\[
F = \frac{U/df_1}{V/df_2}
\sim F(df_1, df_2).
\]

Thus, the F distribution is the ratio of two scaled chi-square random variables.

\subsection{Review of Distributions}

\[
Z \sim N(0,1),
\]
\[
Z^2 \sim \chi^2_1,
\]
\[
t_\nu = \frac{Z}{\sqrt{\chi^2_\nu/\nu}}
\sim t_\nu,
\]
\[
F_{m,n} =
\frac{\chi^2_m/m}{\chi^2_n/n}
\sim F(m,n).
\]

\subsection{Relationship Between $t$ and $F$}

If
\[
t_n = \frac{Z}{\sqrt{\chi^2_n/n}},
\]
then
\[
t_n^2
=
\frac{Z^2}{\chi^2_n/n}
=
\frac{\chi^2_1/1}{\chi^2_n/n}.
\]

Therefore,
\[
t_n^2 \sim F(1,n).
\]

\subsection{Shape of the F Distribution}

\begin{itemize}
\item \textbf{Support:} $F \ge 0$.
\item \textbf{Skewness:} Always right-skewed.
\item \textbf{Degrees of freedom:} Larger degrees of freedom make the distribution less skewed.
\end{itemize}

\textbf{Example:} $F \sim F(5,11)$ produces a right-skewed distribution.

\subsection{Significance of Linear Relationship}

A linear relationship is significant if a large amount of variation in the response is explained by the model.

\begin{itemize}
\item $\SSreg$ should be large relative to $\SST$.
\item $\RSS$ should be small relative to $\SST$.
\item Overall, we check whether $\SSreg$ is significantly larger than $\RSS$.
\end{itemize}

\subsection{F Statistic for Testing Simple Regression}

For testing
\[
H_0:\beta_1 = 0
\qquad
H_1:\beta_1 \neq 0,
\]
the ANOVA decomposition suggests the statistic
\[
F =
\frac{\SSreg/1}{\RSS/(n-2)}.
\]

Under $H_0$,
\[
F \sim F_{1,n-2}.
\]

\subsection{p-Value and Rejection Rule}

\[
\text{p-value} = P(F_{1,n-2} > F_{\text{obs}}).
\]

Reject $H_0$ at level $\alpha$ if
\[
\text{p-value} < \alpha
\]
or equivalently if
\[
F_{\text{obs}} \ge F_{\alpha,1,n-2}.
\]

\textbf{Conclusion:} If $H_0$ is rejected, there is evidence of a statistically significant linear relationship.

\subsection{Mean Squares}

\textbf{Regression mean square:}
\[
MS_{\text{reg}} = \frac{\SSreg}{p}.
\]

\textbf{Residual mean square:}
\[
MS_{\text{res}} = \frac{\RSS}{n-p-1}.
\]

For simple linear regression ($p=1$):
\[
F =
\frac{\SSreg/1}{\RSS/(n-2)}
=
\frac{MS_{\text{reg}}}{MS_{\text{res}}}.
\]

\subsection{ANOVA Table (Simple Linear Regression)}

\begin{center}
\begin{tabular}{cccccc}
\hline
Source & df & Sum of Squares & Mean Square & F \\
\hline
Regression & 1 & $\SSreg$ & $\SSreg/1$ & $\frac{\SSreg/1}{\RSS/(n-2)}$ \\
Residual & $n-2$ & $\RSS$ & $\RSS/(n-2)$ & \\
Total & $n-1$ & $\SST$ & & \\
\hline
\end{tabular}
\end{center}

\subsection{Relationship Between $F$ and $t$ Tests}

For simple linear regression,
\[
F = T^2.
\]

\textbf{Thus, the ANOVA $F$-test and the $t$-test for $\beta_1$ produce the same p-value.}

\pagebreak

\subsection{F-Test in Multiple Linear Regression}

Suppose the model has $p$ predictors:
\[
Y = \beta_0 + \beta_1 X_1 + \beta_2 X_2 + \cdots + \beta_p X_p + \varepsilon.
\]

The overall $F$-test checks whether the predictors collectively explain variation in $Y$.

\subsubsection{Hypothesis Test}

\[
H_0 : \beta_1 = \beta_2 = \cdots = \beta_p = 0
\]
\[
H_1 : \text{at least one } \beta_i \neq 0.
\]

\textbf{Interpretation:} Under $H_0$, none of the predictors help explain the response.

\subsubsection{Test Statistic}

\[
F = \frac{MS_{\text{reg}}}{MS_{\text{res}}},
\]
where
\[
MS_{\text{reg}} = \frac{\SSreg}{p}, \qquad
MS_{\text{res}} = \frac{\RSS}{n - p - 1}.
\]

Thus,
\[
F = \frac{\SSreg/p}{\RSS/(n-p-1)}.
\]

Under $H_0$,
\[
F \sim F_{p,\,n-p-1}.
\]

\subsubsection{p-Value}

\[
p\text{-value} = P\big(F_{p,n-p-1} > F_{\text{obs}}\big),
\]
where $F_{\text{obs}}$ is the observed value of the statistic.

\subsubsection{Rejection Rule}

Reject $H_0$ at level $\alpha$ if
\[
p\text{-value} < \alpha
\]
or equivalently if
\[
F_{\text{obs}} \ge F_{\alpha,p,n-p-1}.
\]

\textbf{Conclusion:} Rejecting $H_0$ implies evidence of a statistically significant linear relationship.

\subsection{ANOVA Table for Multiple Linear Regression}

\begin{center}
\begin{tabular}{cccccc}
\hline
Source & df & Sum of Squares & Mean Square & F \\
\hline
Regression & $p$ & $\SSreg$ & $\SSreg/p$ & $\frac{\SSreg/p}{\RSS/(n-p-1)}$ \\
Residual & $n-p-1$ & $\RSS$ & $\RSS/(n-p-1)$ & \\
Total & $n-1$ & $\SST$ & & \\
\hline
\end{tabular}
\end{center}

\subsection{Error Variance Estimate}

The estimator of the noise variance is
\[
\hat{\sigma}^2 = s^2 = \frac{\RSS}{n-p-1}.
\]

\subsection{F-Test vs t-Test}

\textbf{t-test:} Tests a single parameter,
\[
H_0 : \beta_j = 0.
\]

\textbf{F-test:} Tests joint significance of predictors,
\[
H_0 : \beta_1 = \beta_2 = \cdots = \beta_p = 0.
\]

\subsection{Why Not Run $p$ t-Tests?}

Running $p$ separate t-tests does not test the joint effect of predictors. The $F$-test evaluates whether the predictors together explain significant variation in $Y$.

\subsection{Multicollinearity}

Multicollinearity refers to high correlation among predictor variables. In such cases:
\begin{itemize}
\item Individual $t$-tests may appear insignificant.
\item The overall $F$-test can still be significant.
\end{itemize}

\clearpage
